% Inverting anticryptography --- draft for IACR ePrint
%
% Compiles with pdflatex using only standard packages (no llncs.cls required),
% so it builds on Overleaf out of the box. To submit under the Springer LNCS
% style instead, swap the documentclass for \documentclass{llncs} and move the
% abstract into \begin{abstract}...\end{abstract} inside \maketitle.
\documentclass[11pt,a4paper]{article}

\usepackage[T1]{fontenc}
\usepackage[utf8]{inputenc}
\usepackage{lmodern}
\usepackage{amsmath,amssymb,amsthm}
\usepackage{booktabs}
\usepackage{array}
\usepackage[margin=1in]{geometry}
\usepackage[hidelinks]{hyperref}
\usepackage{microtype}

\newtheorem{proposition}{Proposition}

\title{Inverting anticryptography:\\
a capability-gated interstellar signal\\
as a cipher under brute-force keysearch}
\author{Robert Griffin\\
\small Dxtra Inc.\\
\small \texttt{rtg@dxtra.com}}
\date{}

\begin{document}
\maketitle

\begin{abstract}
The SETI message-design literature is explicitly \emph{anticryptographic}: from
Freudenthal's \emph{Lincos} through Drake's Arecibo tests to contemporary message
design, every construction assumes a sender who maximises decodability---maximum
redundancy, minimum assumed context. We consider the inverse. Suppose a sender,
constrained by transmission cost, wishes the message to be legible \emph{only} to
a receiver above some capability threshold, and therefore makes decoding
deliberately expensive.

We show this construction is a \textbf{symmetric cipher under brute-force
keysearch} with an unusual property: the plaintext is \emph{self-identifying}, so
the receiver needs no crib, and the ciphertext is a \emph{public record} the
sender did not transmit. We then observe that for such a scheme \textbf{Shannon's
unicity distance and the multiple-testing threshold of the receiver's search are
the same quantity}---the sender's difficulty parameter and the receiver's
statistical burden are one number seen from two sides.

Three consequences follow. Required record length scales as $\ln K$, so keyspace
expansion is logarithmically cheap for the sender and linearly expensive for the
receiver---an unusually favourable asymmetry for a gating construction. A
receiver searching below the unicity distance for its keyspace obtains
\textbf{no result}, not a negative one, a distinction we argue must be reported
separately or accumulated coverage estimates are corrupted. And we prove that
linear multi-channel forms cannot carry a signal invisible to pairwise analysis,
which constrains what such a sender could choose.

We report an instantiation over ${\sim}10^4$ keys on 30 public time series, in
which the predicted below-unicity failure occurred and was diagnosed.

\medskip
\noindent\textbf{No hardness result is claimed and no security proof is
transferred.} The contribution is the inversion, the correspondence it makes
visible, and one theorem.

\medskip
\noindent\textbf{Keywords:} anticryptography $\cdot$ SETI $\cdot$ unicity
distance $\cdot$ brute-force keysearch $\cdot$ self-identifying plaintext $\cdot$
time-lock puzzles $\cdot$ multiple hypothesis testing $\cdot$ Kerckhoffs's
principle
\end{abstract}

\section{Introduction}

\subsection{The tradition this inverts}\label{sec:tradition}

Sixty years of interstellar message design share one assumption: \textbf{the
sender wants to be understood.} Freudenthal's \emph{Lincos}~\cite{freudenthal}
builds a language from logic upward so a receiver with no shared context can
bootstrap it. Drake tested an early Arecibo message on colleagues at Green Bank
in 1961 specifically to measure decodability. The tradition has an explicit
name---

\begin{quote}
\textbf{Anticryptography:} the branch of signal synthesis concerned with making
signals easily understandable; ``exactly the opposite goal to that of the science
of cryptography.''~\cite{xenology}
\end{quote}

\noindent---and contemporary work continues it: general-purpose binary languages
for crowdsourced messages, controlled decipherment experiments in which one
researcher encodes and another decodes~\cite{busch}, and public decrypt
challenges~\cite{decrypt}. Every such construction maximises redundancy and
minimises assumed context.

\textbf{This paper considers the inverse.} Suppose a sender is constrained by
transmission cost, and wishes the message to be acted upon \emph{only} by a
receiver above some capability threshold---because a reply from a receiver who
cannot act on it is wasted energy. Such a sender has reason to make decoding
\textbf{deliberately expensive}, and to let the expense do the filtering.

To our knowledge this inversion has not been proposed. The nearest idea in the
literature is the zoo hypothesis~\cite{ball}, which shares the \emph{prediction}
---silence despite presence---but supplies no mechanism, and requires many
parties to abstain indefinitely. A difficulty-gated signal requires nobody to
abstain: it is simply unrecognised until someone spends enough.

\subsection{Why this is a cryptographic object}

Once difficulty is the design goal, the construction becomes a cipher and the
standard analysis applies. Classical cryptanalysis assumes an adversary who must
\emph{recognise} a correct decryption; Shannon fixed the condition under which
that is possible, the \textbf{unicity distance}~\cite{shannon}
\[
  U \;=\; \frac{H(K)}{D},
\]
with $H(K)$ the keyspace entropy and $D$ the plaintext redundancy per unit
length. Below $U$, spurious keys decrypt plausibly and the attack cannot
conclude.

Our setting has three unusual features:
\begin{itemize}
  \item the \textbf{ciphertext is a public record} the sender never
    transmitted---an observational archive a third party collected for unrelated
    reasons;
  \item the \textbf{key is a selection}, not a bit string: which quantities are
    combined and under which functional form;
  \item the \textbf{plaintext is self-identifying}---the receiver knows the key is
    right because a statistic fires. No crib, no known-plaintext pair, no channel
    back to the sender.
\end{itemize}

The receiver's task is a brute-force keysearch with a statistical oracle. The
observation of this paper is that the stopping condition for that search is the
unicity distance in different vocabulary, and that this identity ties the
sender's difficulty parameter to the receiver's statistical burden.

\subsection{Contributions}

\begin{enumerate}
  \item \textbf{The inversion (\S\ref{sec:tradition}).} Sixty years of
    interstellar message design assume a sender maximising decodability. We
    consider a sender minimising it, and show the resulting object is a cipher.
    We find no prior proposal of a difficulty-gated interstellar signal; the zoo
    hypothesis shares the prediction and supplies no mechanism.
  \item \textbf{The correspondence (\S\ref{sec:unicity}).} Unicity distance and
    the multiple-testing threshold are the same quantity---the sender's
    difficulty parameter and the receiver's statistical burden seen from two
    sides. A Bonferroni or \v{S}id\'ak correction over $K$ candidate keys is a
    unicity-distance calculation.
  \item \textbf{The scaling (\S\ref{sec:scaling}).} Required record length grows
    as $\ln K$. Keyspace expansion is logarithmically cheap in data and linearly
    expensive in compute---the opposite of the usual intuition that a larger
    search space demands proportionally more evidence.
  \item \textbf{The void/null distinction (\S\ref{sec:void}).} A search below the
    unicity distance for its keyspace returns \emph{no result}, not a negative
    one. We argue this should be a reported category.
  \item \textbf{Payment inversion (\S\ref{sec:notpow}).} The construction is not
    proof-of-work: the receiver performs the work. We give the consequences for
    what the scheme can and cannot enforce.
  \item \textbf{A structural theorem (\S\ref{sec:theorem}).} For linear
    multi-channel forms, any signal invisible to every pairwise test is invisible
    to the multi-channel form as well. Only genuinely multiplicative forms escape.
  \item \textbf{An empirical instantiation (\S\ref{sec:empirical})} over
    ${\sim}10^4$ keys, in which the predicted below-unicity failure was observed
    and diagnosed.
\end{enumerate}

\section{The construction}\label{sec:construction}

Let $R$ be a public record comprising $n$ observable series
$x_1(t),\dots,x_n(t)$, each of length $N$.

A \textbf{key} is a pair $k=(C,f)$ where $C\subseteq\{1,\dots,n\}$ with $|C|=m$
is a choice of channels and $f\in F$ a functional form combining them. The
keyspace has size
\[
  K \;=\; \binom{n}{m}\,|F|, \qquad H(K)=\log_2 K .
\]

The \textbf{plaintext} is a modulation $s(t)$ imposed on the channels in $C$ such
that $f$ applied to those channels exhibits structure absent under the null
hypothesis for $R$.

The \textbf{oracle} is a statistic $T(k)$ computed from $R$ under key $k$,
together with a null distribution for $T$. The plaintext is self-identifying in
the sense that $T$ exceeds threshold if and only if $k$ is correct---the receiver
needs nothing but $R$ and $T$ to know.

Three properties distinguish this from textbook keysearch.

\paragraph{The ciphertext is not secret and not addressed.} $R$ was collected by
third parties for unrelated purposes, and is public. There is no channel between
the party imposing $s$ and the party searching.

\paragraph{The scheme satisfies Kerckhoffs's principle~\cite{kerckhoffs} in its
strongest form.} Nothing is hidden. The construction, the keyspace, the statistic
and the record may all be published without loss. Whatever difficulty exists lies
entirely in the work of enumeration.

\paragraph{The keyspace is not uniform.} Whoever imposed $s$ had to choose
channels the receiver plausibly measures, on records long enough to hold
structure. This concentrates the key distribution---see \S\ref{sec:weak}.

\section{It is not proof of work: the payment is inverted}\label{sec:notpow}

The construction is often described by analogy to proof-of-work. The analogy
misnames the primitive.

In hashcash~\cite{hashcash} and its descendants the \textbf{sender} performs
work, to price a message or prove commitment. Here the \textbf{receiver} performs
the work, and the sender performs none. The direction of payment is inverted, and
with it the design goal: the scheme does not ration access to the sender's
resources, it conditions legibility on the receiver's.

A more accurate description is a \textbf{proof-of-capability challenge}---a
CAPTCHA inverted. A CAPTCHA is a puzzle a machine is meant to fail. This is a
puzzle only a sufficiently resourced analyst can pass.

The distinction is not terminological. It determines what the scheme can enforce.
A proof-of-work scheme can impose a cost on a sender and verify it. This scheme
can impose nothing on anyone; it can only remain unrecognised until a receiver
spends enough. It is therefore self-enforcing in a weak sense---no party need
cooperate for it to hold---and coercive in no sense at all.

\subsection{Relation to time-lock puzzles}

The nearest technical antecedent is Rivest, Shamir and Wagner's \textbf{time-lock
puzzle}~\cite{rsw}, which requires a chosen amount of computation to open and
whose difficulty is set from \emph{projected future hardware speed}, so that the
puzzle opens at a chosen future date. That is the same reasoning as setting a
keyspace size such that only a receiver above a given compute threshold can
exhaust it.

The constructions differ in one load-bearing respect. A time-lock puzzle is
\textbf{inherently sequential}---its delay cannot be bought off with more
machines, the property later formalised in verifiable delay
functions~\cite{vdf}. \textbf{A keysearch parallelises perfectly.} Our
construction therefore cannot enforce a delay; it can only enforce a total work
requirement. A receiver with $P$ processors finishes in $1/P$ the time.

This makes the scheme strictly weaker than a VDF as a timing primitive, and
strictly more practical as a capability filter: total work is a meaningful
threshold even when elapsed time is not.

\section{Unicity distance is the multiple-testing threshold}\label{sec:unicity}

\subsection{The correspondence}

A receiver enumerating $K$ keys computes $K$ statistics and must decide whether
the largest is evidence of a true key or the expected maximum of $K$ draws from
the null. The standard remedy is a multiplicity correction: test at $\alpha/K$,
or compare against the distribution of $\max T$ over $K$ nulls.

This is Shannon's condition restated. Unicity distance is the ciphertext length
below which spurious keys yield plausible decryptions. The multiple-testing
threshold is the statistic level below which spurious keys fire. \textbf{They are
the same boundary, derived in two vocabularies}---one information-theoretic, one
frequentist.

The practical consequence is that the multiplicity correction in such a search is
not a conservative statistical convention that a cleverer method might avoid. It
is an information-theoretic bound. No test can do better, because below the
unicity distance the information required to distinguish the true key is not
present in the record.

\subsection{Deriving the scaling}\label{sec:derive}

Let $T$ have a null distribution with sub-Gaussian tails. The maximum of $K$
independent draws grows as
\[
  \mathbb{E}\!\left[\max_{1\le i\le K} T_i\right] \;\approx\; \sigma\sqrt{2\ln K}.
\]
A modulation of fractional amplitude $a$ imposed on a record of length $N$ yields
a statistic scaling as $a\sqrt{N}$ against the same $\sigma$. Detection requires
\[
  a\sqrt{N} \;\gtrsim\; \sigma\sqrt{2\ln K}
  \qquad\Longrightarrow\qquad
  \boxed{\,N \;\gtrsim\; \frac{2\sigma^{2}}{a^{2}}\,\ln K\,.}
\]
Since $H(K)=\log_2 K$, this is $N \propto H(K)$---exactly the form of Shannon's
$U = H(K)/D$, with $1/D$ playing the role of the per-sample information about the
modulation.

\subsection{The consequence that should change reporting}\label{sec:void}

If $N < U(K)$, the search cannot conclude. It has not shown the plaintext absent;
it has shown that this record is too short to distinguish $K$ candidate keys, for
any test whatsoever.

We propose that such an outcome be reported as \textbf{void} rather than
\textbf{null}, and that the two be counted separately:

\begin{center}
\begin{tabular}{@{}ll@{}}
\toprule
outcome & meaning \\
\midrule
\textbf{negative} & $N \ge U(K)$; the key is not in the searched space at the stated amplitude \\
\textbf{void} & $N < U(K)$, or no calibrated null exists; the search is uninformative \\
\bottomrule
\end{tabular}
\end{center}

Reporting a void as a negative overstates coverage and, in a setting where
searches accumulate, corrupts any estimate of how much of the keyspace has been
eliminated.

\section{Keyspace expansion is cheap in data, expensive in compute}\label{sec:scaling}

\subsection{The asymmetry}

From \S\ref{sec:derive}, enlarging the keyspace by a factor of ten costs a factor
of $\ln(10K)/\ln K$ in record length---for $K\sim10^4$, about 25\%. It costs a
factor of ten in computation.

This inverts a common intuition. Analysts frequently restrict a search space to
``preserve statistical power,'' reasoning that fewer tests means a laxer
threshold. The scaling says the saving is logarithmic and the loss is the risk of
excluding the true key entirely. \textbf{Where the record is fixed and compute is
elastic---the usual modern situation---the correct move is to search the larger
space.}

Two caveats attach. The ladder below counts \emph{nominal} keys; the
\textbf{reachable} keyspace may be smaller, since a key whose channels share no
overlapping record cannot be tested---in the instantiation of
\S\ref{sec:empirical} the reachable pair space was 88\% of nominal and the
reachable triple space 72\%. The correct $K$ for a threshold calculation is the
number of keys actually tested, not the number enumerable.

\begin{center}
\begin{tabular}{@{}lrrr@{}}
\toprule
tier & keys $K$ & $H(K)$ bits & record length, relative \\
\midrule
pairs                      & $4.35\times10^{2}$ & 10.8 & $1.00\times$ \\
triples                    & $4.06\times10^{3}$ & 14.0 & $1.30\times$ \\
quadruples                 & $2.74\times10^{4}$ & 16.7 & $1.56\times$ \\
triples over 60, 10 forms  & $3.42\times10^{4}$ & 18.4 & $1.71\times$ \\
\bottomrule
\end{tabular}
\end{center}

A sixty-fold keyspace expansion costs 71\% more record.

\subsection{Weak keys, and why brute force is the wrong first attack}\label{sec:weak}

The keyspace is not uniform. Whoever chose $C$ was constrained to channels the
receiver plausibly measures, at cadences and durations sufficient to carry
structure. The key distribution concentrates on the best-instrumented and most
obvious observables.

In cryptographic terms these are \textbf{weak keys}, and the standard response
applies: \emph{no password cracker begins with brute force; it begins with a
wordlist.} An ordered search over a prior-weighted enumeration reaches a probable
key far sooner in expectation, at no cost in coverage provided the full
enumeration is eventually completed.

We stress the limit of this argument. The prior is the receiver's reasoning about
a sender's choices and is not testable. It justifies search \textbf{order}, which
costs nothing if wrong. It does not justify restricting the search
\textbf{space}.

\section{A structural constraint on multi-channel carriers}\label{sec:theorem}

A natural question is whether a modulation can be hidden from all pairwise
analysis while remaining visible to a higher-order form. For linear forms the
answer is no.

\begin{proposition}
Let a modulation $s(t)$ be imposed with coefficients $c_k$ on channel $k$.
Consider the linear four-channel form with coefficients $(1,-3,3,-1)$---the third
finite difference---and the pairwise forms $x_i-x_j$. If the modulation is
invisible to every pairwise form, it is invisible to the four-channel form.
\end{proposition}

\begin{proof}
The pairwise form on $(i,j)$ picks up $(c_i-c_j)s$. Invisibility to every pair
requires $c_i=c_j$ for all $i,j$, hence all $c_k$ equal to some $c$. The
four-channel form then picks up $(c-3c+3c-c)s=0$.
\end{proof}

The result extends to any linear form whose coefficients sum to zero, which
includes the finite differences and cross-differences commonly used as
higher-order statistics. \textbf{Such forms are not independent carriers}; they
are differences of pairwise information.

Escaping the proposition requires a genuinely multiplicative form. A four-way
phase alignment, living in the trispectrum, leaves every pairwise
cross-correlation untouched and is therefore a real higher-order carrier.

The practical import for search design is that enumerating linear higher-order
forms enlarges $K$---and by \S\ref{sec:derive} raises the detection threshold for
every key---while adding no reachable plaintext. It is a strictly bad trade.

\section{Empirical instantiation}\label{sec:empirical}

We instantiated the construction on 30 public observational time series,
enumerating pairs and triples under four functional forms~\cite{griffin}.

The reachable keyspace is smaller than the nominal one, and the distinction
matters for \S\ref{sec:derive}: a key whose channels have no overlapping record
is not a key the receiver can test at all. Of 435 nominal pairs, 382 had the
required 2{,}000-sample overlap, giving \textbf{1{,}074 completed tests} against
60{,}000 circular shifts each. Of 4{,}060 nominal triples, 2{,}934 were reachable
and 2{,}749 completed, giving \textbf{10{,}996 tests}---117 million statistic
evaluations in 26 hours on one 88-core machine.

The pair search completed and returned negative at a stated amplitude.
\textbf{The triple search returned void, in the precise sense of
\S\ref{sec:void}}, and the diagnosis is instructive: no calibrated null
distribution could be constructed for the higher-order statistics. Three
candidate nulls were built and all three failed, each because it required a false
modelling assumption about the underlying process:

\begin{center}
\begin{tabular}{@{}p{3.1cm}p{4.6cm}p{5.4cm}@{}}
\toprule
null & assumption & failure \\
\midrule
circular shift & pairwise structure irrelevant to a three-body statistic
  & destroys it; ninefold tail excess \\
common phase screen & series are stationary
  & they are not; real data at median $z=45$ \\
envelope-preserving & the nonlinearity is pure amplitude modulation
  & absorbs the signal; $4.25\times$ over-firing \\
\bottomrule
\end{tabular}
\end{center}

Additionally, and in agreement with \S\ref{sec:theorem}, the two linear
higher-order forms were swamped at $z\approx50$ while the single multiplicative
form sat near its null at median $z=5.1$---the predicted signature of forms that
are not independent carriers.

The full search, its code and its data are public~\cite{griffin}.

\section{What is not claimed}

\begin{itemize}
  \item \textbf{No hardness result.} The mapping to a cipher is structural. No
    security proof, reduction or hardness assumption is transferred, and none is
    implied. The keyspace arithmetic and the $\sqrt{2\ln K}$ threshold scaling
    stand on their own.
  \item \textbf{No claim of novelty for the components.} Unicity distance is
    Shannon's. Self-identifying plaintext is the standing assumption beneath
    every brute-force attack. Time-lock puzzles are RSW's. The claimed
    contribution is the inversion of \S\ref{sec:tradition}, the correspondence of
    \S\ref{sec:unicity}, the void/null distinction, and the proposition of
    \S\ref{sec:theorem}.
  \item \textbf{The weak-key prior is not testable.} It orders a search; it does
    not bound one.
  \item \textbf{The empirical instantiation detected nothing.} It is reported as
    an existence proof for the failure mode of \S\ref{sec:void}, not as evidence
    for any hypothesis about its subject matter.
\end{itemize}

\section{Open questions}

\begin{enumerate}
  \item Does the correspondence of \S\ref{sec:unicity} extend to keyspaces with
    non-trivial structure---where keys are related, so that the $K$ tests are
    dependent? The sub-Gaussian maximum is then an over-estimate and an effective
    $K$ smaller than the nominal one may apply.
  \item Is there a construction with self-identifying plaintext that is also
    inherently sequential, recovering the VDF property lost in
    \S\ref{sec:notpow}?
  \item Can the void/null distinction be given a sharper operational test than
    ``no calibrated null exists''---ideally one computable before a search is run?
\end{enumerate}

\begin{thebibliography}{99}

\bibitem{griffin} Griffin, R.: A combination-space search for embedded
technosignatures in solar and heliospheric archives (2026). Code and data:
\url{https://github.com/dxtrainc/behavioral-seti}

\bibitem{shannon} Shannon, C.E.: Communication theory of secrecy systems. Bell
System Technical Journal \textbf{28}(4), 656--715 (1949)

\bibitem{rsw} Rivest, R.L., Shamir, A., Wagner, D.A.: Time-lock puzzles and
timed-release crypto. MIT/LCS/TR-684 (1996)

\bibitem{vdf} Boneh, D., Bonneau, J., B\"unz, B., Fisch, B.: Verifiable delay
functions. In: CRYPTO 2018

\bibitem{hashcash} Back, A.: Hashcash --- a denial of service counter-measure
(2002)

\bibitem{kerckhoffs} Kerckhoffs, A.: La cryptographie militaire. Journal des
sciences militaires \textbf{IX}, 5--38 (1883)

\bibitem{xenology} Freitas, R.A.: Xenology, \S24.2.4 --- on anticryptography as
the inverse of cryptography in signal synthesis. \emph{[verify]}

\bibitem{busch} Busch, M.W., Reddick, R.M.: Testing SETI message designs.
arXiv:0911.3976 (2009). \emph{[verify]}

\bibitem{freudenthal} Freudenthal, H.: Lincos: Design of a Language for Cosmic
Intercourse. North-Holland (1960)

\bibitem{decrypt} Heller, R., Pudritz, R.E.: Decryption of messages from
extraterrestrial intelligence using the power of social media --- the SETI
Decrypt Challenge. arXiv:1706.00653 (2017). \emph{[verify --- authorship not
confirmed]}

\bibitem{ball} Ball, J.A.: The zoo hypothesis. Icarus \textbf{19}, 347 (1973)

\end{thebibliography}

\end{document}
